\scanpage{6}\mspage{3.}
3. \uline{The} ``\uline{minimal}'' \uline{Weierstrass} \uline{equations} \uline{over} \uline{a} \uline{valuation} \uline{ring}. Let $v$ be a discrete valuation of $K$, with valuation ring $R$, prime ideal $\mathfrak{p}$, and residue field $k=R/\mathfrak{p}$. Let $E$ be an elliptic curve over $K$, with a $K\cdot$rational point $O$.

\uline{Defn 3.1}: An equation for $E$ of the form (1.1) is minimal (with respect to $v$) if $v(a_i)\geqslant0$ for all $i$ and if $v(\Delta)$ is minimal, subject to that condition.

\uline{Theorem 3.2}\unsure{.|:} \uline{A} \uline{minimal} \uline{equation} \inserted{for $E$} \uline{exists}, \uline{and} \uline{is} \uline{unique} \uline{up} \uline{to} \uline{a} \uline{change} \uline{of} \uline{coördinates} \uline{of} \uline{the} \uline{form} (2.1) \uline{with} $r,s,t\in R$ \uline{and} $u$ \uline{invertible} \uline{in} $R$.

Existence is obvious. Let $y^{\prime\,2}+a_1'x'y'+\unsure{\ldotp\,\ldotp|\ldots}$ and $y^2+a_1xy+\ldots$ be two minimal equations for the same $E$. Since $\Delta\ne0$ and $v(\Delta')=v(\Delta)$ we conclude from\deleted{\illegible{1}} (2.3) that $v(u)=0$. Now from the transformation of $b_8$ and $b_6$ in (2.2) we see that $3r\in R$ and $4r\in R$, hence $r\in R$. \inserted{Next} The transformation of $a_2$ shows that $s\in R$, and that of $a_6$ shows that $t\in R$.

\uline{Corollary 3.3:} \uline{The} \uline{differential} $\omega$ \uline{associated} \uline{with} \uline{a} \uline{minimal} \deleted{Weierstrass form} \uline{is} \uline{unique} \uline{up} \uline{to} \uline{a} \uline{unit} \uline{in} $R$.

\uline{Remarks}: 1) If $a_i\in R$ and $v(\Delta)<12$, then the equation (1.1) is automatically minimal. The converse is true if $j\in R$ and $p=\operatorname{char}k\ne2,3$. A complete algorithm for reducing to minimal form in all \unsure{cases|\deleted{c}cases} is given below.

2) Let $\mathscr{M}$ be a collection of discrete valuations of $K$. For each $v\in\mathscr{M}$, let $y_v^2+a_{1,v}x_vy_v+\ldotp\,\ldotp\,\ldotp\,\ldotp$ be a minimal equation for $E$ relative to $v$, with corresponding differential $\omega_v$ and discriminant $\Delta_v$. Then the ``divisor'' $\displaystyle \mathfrak{d}_E=\sum_{v\in\mathscr{M}}v(\Delta_v)v$ should be regarded as \uline{the} \uline{discriminant} \uline{of} $E$.

Let $F(x,y)=y^2+a_1xy+\ldots$ be an arbitrary Weierstrass equation for $E$ over $K$, and let $\omega_F$ and $\Delta_F$ be the corresponding differential and discriminant. Then the \uline{class} of the divisor $\displaystyle \mathfrak{u}_F=\sum_v v\left(\frac{\omega}{\omega_v}\right)v$ is independent of $F$, and we have $12\mathfrak{u}_F=\mathfrak{d}_E-(\Delta_F)$, so that $12\mathfrak{u}_F\sim\mathfrak{d}_E$.

3) In case $\mathscr{M}$ is the set of valuations associated with a principal ideal domain $D$ with field of fractions $K$, then it is easy to see that we can find one equation $F$ which simultaneously minimal for all $v$, so that $\mathfrak{u}_F=0$, and $\mathfrak{d}=(\Delta_F)$.

4) If $j\in R$, then $\Delta$ divides $c_4^3$ and $c_6^2$. From (1.7) we see that if $48\mathfrak{p}^4\mid c_4$ and $864\mathfrak{p}^6\mid c_6$, then the equation is not minimal. Hence \uline{if} $j\in R$ \uline{and} \uline{equation} \uline{minimal} \uline{we} \uline{have} $v(\Delta)<$ \deleted{\illegible{?}} $12+12v(2)+6v(3)$

% archive sheet number: 5
